%==============================================================================
% Paper B — JUTI ITS (SINTA 3)
% FedX-Palm: aplikasi FL + Grad-CAM++ untuk deteksi kematangan sawit.
%
% CARA PAKAI DI OVERLEAF:
%   1. New Project -> Upload Project -> unggah folder ini (main.tex + figures/).
%   2. Compiler: pdfLaTeX. Klik Recompile.
%   3. WAJIB ganti SEMUA \tbd{...} (muncul MERAH) dengan angka NYATA dari
%      aggregate_results.py & evaluate_xai.py SEBELUM submit. Jangan upload
%      kalau masih ada teks merah.
%   4. Idealnya, pindahkan isi ke template resmi JUTI bila tersedia; berkas ini
%      adalah manuskrip lengkap berformat jurnal dua-kolom yang setara.
%==============================================================================
\documentclass[10pt,twocolumn]{article}

\usepackage[T1]{fontenc}
\usepackage[utf8]{inputenc}
\usepackage[bahasa,english]{babel}
\usepackage{graphicx}
\usepackage{booktabs}
\usepackage{amsmath}
\usepackage{amssymb}
\usepackage{authblk}
\usepackage[hidelinks]{hyperref}
\usepackage{xcolor}
\usepackage{geometry}
\geometry{a4paper,margin=2.2cm}

% Placeholder hasil — tampil MERAH agar tidak terlewat sebelum submit.
\newcommand{\tbd}[1]{\textcolor{red}{\textbf{[ISI: #1]}}}

\title{\textbf{Deteksi Kematangan Tandan Buah Segar Kelapa Sawit Enam Kelas
Berbasis \textit{Federated Learning} dengan Penjelasan Visual Grad-CAM++
Tervalidasi untuk Pendukung Keputusan Panen}}

\author[1]{Nama Penulis Pertama}
\author[2]{Nama Pembimbing}
\affil[1]{Program Studi ..., Fakultas ..., Universitas ...; \texttt{email@domain}}
\affil[2]{Departemen ..., Institusi ...; \texttt{email@domain}}

\date{}

\begin{document}
\maketitle

%------------------------------------------------------------------ ABSTRAK ID
\selectlanguage{bahasa}
\begin{abstract}
\noindent
Penentuan waktu panen tandan buah segar (TBS) kelapa sawit yang tepat sangat
menentukan rendemen dan mutu minyak, namun penilaian manual bersifat subjektif
dan tidak konsisten. Otomasi berbasis visi komputer menjanjikan keseragaman,
tetapi pengumpulan citra perkebunan secara terpusat terbentur kerahasiaan data
operasional. Penelitian ini mengevaluasi apakah pendekatan \textit{Federated
Learning} (FL)---yang melatih model bersama tanpa memindahkan citra mentah
antar-kebun---dapat menghasilkan deteksi kematangan TBS enam kelas yang akurat
sekaligus dapat dijelaskan, sehingga layak menjadi pendukung keputusan panen.
Detektor YOLOv11n dilatih secara federasi (FedAvg) pada data Non-IID (partisi
Dirichlet) hasil pembagian anti-kebocoran berbasis identitas tandan
(\textit{bunch\_id}), lalu kualitas penjelasannya dievaluasi menggunakan
Grad-CAM++ melalui metrik \textit{Average Drop} (AD) dan \textit{Focus
Retention Rate} (FRR) per-kelas. Model federasi mencapai mAP@0,5 =
\tbd{B2 mAP50} (pembanding tersentral 0,787), dengan AD = \tbd{AD global} dan
FRR = \tbd{FRR global}. Analisis per-kelas menunjukkan \tbd{kelas termudah/
tersulit dijelaskan}. Hasil ini mengindikasikan FL dapat menyediakan deteksi
kematangan sawit yang menjaga privasi data antar-kebun tanpa mengorbankan
keterjelasan keputusan model.

\noindent\textbf{Kata kunci:} kelapa sawit, deteksi kematangan, \textit{federated
learning}, YOLOv11, \textit{explainable AI}, Grad-CAM++, pertanian presisi.
\end{abstract}

%------------------------------------------------------------------ ABSTRACT EN
\selectlanguage{english}
\begin{abstract}
\noindent\textit{Timely harvesting of oil palm fresh fruit bunches (FFB) is
decisive for oil yield and quality, yet manual ripeness assessment is
subjective. Computer-vision automation promises consistency, but centrally
aggregating plantation imagery conflicts with operational-data
confidentiality. This study evaluates whether Federated Learning (FL)---
training a shared model without moving raw imagery across estates---can deliver
six-class FFB ripeness detection that is both accurate and explainable. A
YOLOv11n detector is trained federatively (FedAvg) on Non-IID (Dirichlet)
partitions from a leakage-free bunch-identity split, and its explanation
quality is assessed with Grad-CAM++ via Average Drop (AD) and Focus Retention
Rate (FRR) per class. The federated model attains mAP@0.5 = }\tbd{B2 mAP50}
\textit{ (centralized reference 0.787), with AD = }\tbd{AD}\textit{ and FRR = }
\tbd{FRR}\textit{. The results indicate FL can provide privacy-preserving palm
ripeness detection without sacrificing decision explainability.}

\noindent\textit{\textbf{Keywords:} oil palm, ripeness detection, federated
learning, YOLOv11, explainable AI, Grad-CAM++, precision agriculture.}
\end{abstract}

\selectlanguage{bahasa}

%==============================================================================
\section{Pendahuluan}
%==============================================================================
Industri kelapa sawit merupakan tulang punggung agribisnis Indonesia dan
Malaysia yang bersama-sama menyumbang lebih dari 85\% produksi minyak sawit
dunia~\cite{usda2024}. Pada tingkat operasional, keputusan paling menentukan
rendemen adalah ketepatan waktu panen TBS: tandan yang dipanen terlalu dini
berkadar minyak rendah, sedangkan yang terlambat meningkatkan kadar Asam Lemak
Bebas dan menurunkan mutu \textit{Crude Palm Oil}~\cite{junos2022}. Penilaian
manual oleh mandor bersifat subjektif dengan tingkat kesalahan 15--25\% pada
pencahayaan suboptimal~\cite{sabri2017}, sehingga otomasi berbasis visi
komputer menjadi kebutuhan nyata.

Detektor objek modern seperti YOLOv11 mampu melokalisasi dan mengklasifikasi
beberapa tandan dalam satu bingkai secara \textit{real-time}~\cite{yolov11},
namun pelatihannya memerlukan banyak citra beragam. Pengumpulan terpusat
terbentur kenyataan bahwa citra perkebunan adalah informasi bisnis sensitif
yang enggan dibagikan antar-perusahaan. \textit{Federated Learning}
(FL)~\cite{fedavg} menjawab kendala ini dengan melatih model bersama tanpa
memindahkan data mentah; dalam taksonomi~\cite{kairouz2021}, skenario
perkebunan termasuk \textit{cross-silo}.

Namun, akurasi saja tidak cukup agar model dipercaya operator. Sistem pendukung
keputusan pertanian memerlukan keterjelasan (\textit{explainability}): operator
perlu memahami \emph{mengapa} model mengklasifikasikan suatu tandan. Metode
\textit{Explainable AI} (XAI) seperti Grad-CAM++~\cite{gradcampp} menghasilkan
peta panas yang menyoroti wilayah penentu prediksi, dan kualitasnya dapat
diukur secara kuantitatif.

Berbeda dengan studi pendahuluan penulis yang berfokus pada perilaku
\textit{Differential Privacy} pada FL~\cite{paperA}, penelitian ini berfokus
pada pertanyaan aplikatif: \emph{apakah deteksi kematangan TBS berbasis FL
dapat sekaligus akurat dan dapat dijelaskan dengan andal sehingga layak menjadi
pendukung keputusan panen?} Kontribusi penelitian ini: (1) evaluasi
keterjelasan mendalam deteksi kematangan sawit federasi melalui analisis
Grad-CAM++ \emph{per-kelas}; (2) protokol pembagian data anti-kebocoran
berbasis \textit{bunch\_id}; dan (3) bukti empiris bahwa federasi tidak
menurunkan kualitas penjelasan secara signifikan dibanding pelatihan tersentral.

%==============================================================================
\section{Tinjauan Pustaka}
%==============================================================================
\subsection{Deteksi Kematangan Kelapa Sawit}
\tbd{1 paragraf: ringkas 3--4 studi deteksi/klasifikasi kematangan sawit
terdahulu; tekankan mayoritas mengasumsikan pelatihan terpusat dan tidak
mengevaluasi keterjelasan secara kuantitatif---celah yang diisi paper ini.}

\subsection{Federated Learning}
FedAvg~\cite{fedavg} mengagregasi pembaruan model lokal terbobot ukuran data;
skenario \textit{cross-silo}~\cite{kairouz2021} cocok untuk konsorsium
perkebunan. Heterogenitas antar-kebun disimulasikan dengan partisi
Dirichlet~\cite{hsu2019}.

\subsection{Explainable AI dan Metrik Faithfulness}
Grad-CAM~\cite{gradcam} dan penyempurnaannya Grad-CAM++~\cite{gradcampp}, yang
lebih akurat saat terdapat banyak instans objek sejenis---kondisi lazim pada
citra TBS. Kualitas penjelasan diukur via \textit{Average Drop} dan \textit{Focus
Retention Rate}.

%==============================================================================
\section{Metodologi}
%==============================================================================
\subsection{Dataset dan Pembagian Anti-Kebocoran}
Dataset terdiri atas 10.814 citra dari 89 tandan unik dengan anotasi enam kelas:
\textit{Abnormal}, \textit{Empty Bunch}, \textit{Overripe}, \textit{Ripe},
\textit{Underripe}, \textit{Unripe}. Karena beberapa citra berasal dari tandan
yang sama, pembagian acak biasa berisiko membocorkan tandan ke \textit{train}
dan \textit{test} sekaligus. Penelitian ini menerapkan pembagian terstratifikasi
berbasis \textit{bunch\_id} pada tingkat tandan (Tabel~\ref{tab:split}), dengan
audit memastikan nol irisan tandan antar-split.

\begin{table}[t]
\centering
\caption{Pembagian data anti-kebocoran berbasis \textit{bunch\_id}.}
\label{tab:split}
\begin{tabular}{lrr}
\toprule
Split & Citra & Tandan \\
\midrule
Train & 9.094 & 72 \\
Valid & 769 & 8 \\
Test  & 951 & 9 \\
\bottomrule
\end{tabular}
\end{table}

\subsection{Partisi Federasi dan Pelatihan}
Data \textit{train} dipartisi ke \tbd{K, mis. 4} klien dengan distribusi
Dirichlet~\cite{hsu2019}. Detektor YOLOv11n (\textasciitilde2,6 juta parameter)
diinisialisasi dari bobot pra-latih COCO; seluruh BatchNorm dikonversi menjadi
GroupNorm agar konsisten dengan jalur eksperimen privasi penulis~\cite{paperA}.
Pelatihan federasi memakai FedAvg, lima ronde komunikasi, dua epoch lokal per
ronde (SGD, \textit{lr}~0,01, batch~16, citra~640$\times$640). Pembanding
tersentral dilatih 50 epoch. Aspek \textit{Differential Privacy} ditangani dan
dianalisis pada publikasi pendahuluan penulis~\cite{paperA}, sehingga tidak
diulang di sini.

\subsection{Evaluasi Keterjelasan}
Grad-CAM++ diterapkan pada lapisan konvolusi terakhir sebelum kepala deteksi.
Pada $N=$\tbd{jumlah} citra uji dihitung, per-kelas dan global:
\textit{Average Drop} (AD; makin rendah makin baik) dan \textit{Focus Retention
Rate} (FRR; proporsi intensitas \textit{heatmap} di dalam \textit{bounding box}
buah, sekaligus proksi kesesuaian agronomis).

%==============================================================================
\section{Hasil dan Pembahasan}
%==============================================================================
\subsection{Akurasi Deteksi}
Tabel~\ref{tab:acc} menyajikan akurasi model federasi dibanding pembanding
tersentral. \tbd{1 kalimat: berapa selisih/biaya FL, dan apakah masih berguna.}

\begin{table}[t]
\centering
\caption{Akurasi deteksi pada himpunan uji.}
\label{tab:acc}
\begin{tabular}{lrrrr}
\toprule
Model & mAP@0,5 & mAP@0,5:0,95 & P & R \\
\midrule
Tersentral (acuan) & 0,787 & 0,672 & 0,815 & 0,833 \\
Federasi (FedAvg) & \tbd{} & \tbd{} & \tbd{} & \tbd{} \\
\bottomrule
\end{tabular}
\end{table}

\subsection{Kualitas Penjelasan Per-Kelas (Kontribusi Inti)}
Tabel~\ref{tab:xai} memuat AD dan FRR per-kelas kematangan.
\tbd{2 paragraf: kelas mana paling andal dijelaskan (mis. Ripe/Overripe karena
isyarat warna kuat) vs paling lemah (mis. Abnormal/Empty Bunch); kaitkan dengan
isyarat agronomis warna kulit, brondolan, tekstur; nilai aplikatif: operator
lebih waspada pada kelas dengan penjelasan lemah.}

\begin{table}[t]
\centering
\caption{Faithfulness Grad-CAM++ per-kelas (AD: \%, makin rendah makin baik;
FRR: 0--1, makin tinggi makin baik).}
\label{tab:xai}
\begin{tabular}{lrr}
\toprule
Kelas & AD (\%) & FRR \\
\midrule
Unripe & \tbd{} & \tbd{} \\
Underripe & \tbd{} & \tbd{} \\
Ripe & \tbd{} & \tbd{} \\
Overripe & \tbd{} & \tbd{} \\
Empty Bunch & \tbd{} & \tbd{} \\
Abnormal & \tbd{} & \tbd{} \\
\midrule
\textbf{Global} & \tbd{} & \tbd{} \\
\bottomrule
\end{tabular}
\end{table}

\subsection{Pengaruh Federasi terhadap Keterjelasan}
\tbd{1 paragraf: bandingkan AD/FRR federasi vs tersentral; jika selisih kecil
maka federasi tidak merusak keterjelasan---temuan utama yang memperkuat
kelayakan FL untuk pendukung keputusan.}

\subsection{Analisis Kualitatif}
Gambar~\ref{fig:xai} menampilkan \textit{heatmap} Grad-CAM++ representatif per
kelas. \tbd{diskusi: sorotan jatuh pada permukaan buah/brondolan (benar) atau
latar (salah).}

\begin{figure}[t]
\centering
% Ganti dengan figures/xai_per_class.png hasil evaluate_xai / pembuatan figur.
\fbox{\parbox[c][3.2cm][c]{0.9\linewidth}{\centering \tbd{sisipkan
figures/xai\_per\_class.png}}}
\caption{Penjelasan Grad-CAM++ per kelas kematangan TBS.}
\label{fig:xai}
\end{figure}

%==============================================================================
\section{Kesimpulan}
%==============================================================================
\tbd{1 paragraf} Penelitian ini menunjukkan deteksi kematangan TBS enam kelas
berbasis \textit{Federated Learning} dapat mencapai akurasi \tbd{kategori}
sekaligus mempertahankan keterjelasan yang \tbd{tinggi/memadai} menurut metrik
\textit{faithfulness} Grad-CAM++. Analisis per-kelas mengungkap bahwa \tbd{}.
Federasi \tbd{tidak/sedikit} menurunkan kualitas penjelasan dibanding pelatihan
tersentral, sehingga FL layak menjadi fondasi pendukung keputusan panen yang
menjaga privasi data antar-kebun. Penelitian lanjut mencakup integrasi DP-SGD
penuh~\cite{paperA}, validasi multi-kebun nyata, dan \textit{deployment} edge.

%==============================================================================
\section*{Ucapan Terima Kasih}
\tbd{opsional}

%==============================================================================
\begin{thebibliography}{99}
\bibitem{usda2024} USDA FAS, \textit{Oilseeds: World Markets and Trade}, 2024.
\bibitem{junos2022} M. H. Junos \textit{et al.}, ``Automatic detection of oil
palm fruits from UAV images using an improved YOLO model,'' \textit{The Visual
Computer}, vol.~38, 2022.
\bibitem{sabri2017} N. Sabri \textit{et al.}, ``Palm oil FFB ripeness grading
using color features,'' \textit{J. Fundam. Appl. Sci.}, 2017.
\bibitem{yolov11} R. Khanam and M. Hussain, ``YOLOv11: An Overview of the Key
Architectural Enhancements,'' arXiv:2410.17725, 2024.
\bibitem{fedavg} B. McMahan \textit{et al.}, ``Communication-Efficient Learning
of Deep Networks from Decentralized Data,'' in \textit{AISTATS}, 2017.
\bibitem{kairouz2021} P. Kairouz \textit{et al.}, ``Advances and Open Problems
in Federated Learning,'' \textit{Found. Trends Mach. Learn.}, 2021.
\bibitem{hsu2019} T.-M. H. Hsu \textit{et al.}, ``Measuring the Effects of
Non-Identical Data Distribution for Federated Visual Classification,''
arXiv:1909.06335, 2019.
\bibitem{gradcam} R. R. Selvaraju \textit{et al.}, ``Grad-CAM: Visual
Explanations from Deep Networks via Gradient-based Localization,'' in
\textit{ICCV}, 2017.
\bibitem{gradcampp} A. Chattopadhyay \textit{et al.}, ``Grad-CAM++: Generalized
Gradient-Based Visual Explanations for Deep CNNs,'' in \textit{WACV}, 2018.
\bibitem{paperA} [Penulis], ``FedX-Palm: Per-Sample DP-SGD versus Client-Level
DP-FedAvg for Privacy-Preserving Oil Palm Ripeness Detection on YOLOv11,''
\tbd{venue Scopus Q3 / status}.
\end{thebibliography}

\end{document}
